% Lettre de motivation - Stephane KPOVIESSI
% Candidature a la Maitrise en gestion (M. Sc.) - strategie - HEC Montreal - projet supervise
% Repond aux quatre questions officielles de la liste des pieces requises.
% Mise en forme demandee : Arial 12, interligne 1,5 (helvet = Nimbus Sans, metriquement compatible).
% Une reponse par page : la contrainte "une page maximum par reponse" se verifie a l'oeil.
% Compile target: pdfLaTeX

\documentclass[12pt,a4paper]{article}

\usepackage[a4paper,top=1.9cm,bottom=1.9cm,left=2.4cm,right=2.4cm]{geometry}
\usepackage[utf8]{inputenc}
\usepackage[T1]{fontenc}
\usepackage[scaled=1.0]{helvet}
\renewcommand{\familydefault}{\sfdefault}
\usepackage{xcolor}
\usepackage[hidelinks]{hyperref}
\usepackage{titlesec}

\definecolor{accent}{HTML}{1F4E79}
\definecolor{muted}{HTML}{5A6165}
\definecolor{rule}{HTML}{D5E2EC}

% Interligne 1,5 : \linespread vient du noyau LaTeX (setspace n'est pas installe).
\linespread{1.5}
\setlength{\parindent}{0pt}
\setlength{\parskip}{0.3em}
\pagestyle{plain}

\titleformat{\section}
  {\normalsize\bfseries\color{accent}}
  {}{0pt}{}
\titlespacing*{\section}{0pt}{0.9em}{0.15em}

\begin{document}
\thispagestyle{empty}

{\LARGE\bfseries Stéphane KPOVIESSI}\par
\vspace{0.25em}
{\footnotesize\color{muted}Cergy / Île-de-France \textbullet{} +33 7 48 59 47 13 \textbullet{} oastephanekpoviessi@gmail.com\\
linkedin.com/in/stephanekpoviessi \textbullet{} seshat-ai.com}\par
\vspace{0.25em}
{\color{rule}\hrule height 1pt}
\vspace{0.3em}

{\bfseries Lettre de motivation \textbullet{} Maîtrise en gestion (M. Sc.), stratégie, projet supervisé \textbullet{} HEC Montréal \textbullet{} automne 2027}\par

\section{1. Choix du programme d'études}
{\itshape\footnotesize Quelles préoccupations professionnelles et intellectuelles vous motivent à entreprendre une maîtrise en gestion? Expliquez votre choix de spécialisation.}\par
\vspace{0.2em}

Mes études d'ingénieur m'ont appris à produire des analyses fiables. Mes premières expériences professionnelles m'ont montré qu'au-delà de leur exactitude, leur véritable valeur réside dans leur capacité à éclairer les décisions. Chez Allianz France, six mois de fiabilisation d'indicateurs de pilotage : pour sécuriser un indicateur, il faut d'abord s'entendre sur ce qu'il mesure et sur ce qu'il engage. Chez Lyne \& Frères, deux ans d'analyses et de tableaux de bord pour la direction : une analyse ne vaut que si quelqu'un s'en saisit pour choisir entre deux options.

Je me suis arrêté à ce point : je sais produire la mesure, je ne sais pas encore conduire la décision qu'elle éclaire. Les arbitrages (priorités, ressources, positionnement d'une offre) se jouent dans des conversations où je veux prendre ma place. C'est pourquoi je vise une maîtrise en gestion plutôt qu'une spécialisation fonctionnelle : mon objet n'est pas de devenir le spécialiste d'une fonction, mais de comprendre une entreprise dans son ensemble et de peser sur ses décisions.

La spécialisation stratégie répond à cette limite : identification des enjeux et des options, formulation de recommandations et de plans d'action, évaluation des risques et des occasions. Ce sont les trois gestes qui manquent à mon parcours. Son caractère interdisciplinaire compte autant : comportement organisationnel, consultation, management, éthique et entrepreneuriat complètent une formation d'ingénieur et m'apportent les regards qui manquent à un profil technique.

Je choisis le cheminement avec projet supervisé plutôt que le mémoire, parce que mon objectif est professionnel : travailler sur un problème réel, avec ses contraintes et ses parties prenantes, et rendre des recommandations utilisables. Mon diplôme d'ingénieur relève d'un domaine que vous considérez connexe : j'apporte l'analyse et le traitement des données, la maîtrise m'apportera la gestion.

\newpage
\section{2. Plan de carrière}
{\sloppy\itshape\footnotesize Quel genre d'emploi pensez-vous occuper après votre maîtrise ès sciences en gestion (ex. type d'entreprise, postes et responsabilités visés)?\par}
\vspace{0.35em}

À court terme, je cherche un poste où l'analyse sert directement la décision, dans une direction stratégie, transformation ou performance, ou en conseil : analyste en intelligence d'affaires, conseiller en gestion ou consultant en cabinet. J'y veux des responsabilités concrètes : comprendre une situation, construire un diagnostic, proposer des options argumentées et les défendre devant ceux qui décident. Le secteur m'importe moins que la nature du travail. Une PME ou une ETI en croissance me conviendrait autant qu'un grand groupe : j'ai travaillé dans une grande organisation comme dans une petite entreprise, et les deux exigent de s'adapter à des moyens et à des rythmes différents.

À moyen terme, je veux prendre la responsabilité de projets de transformation, puis de mandats (coordination de projets, gestion de mandat, conseil), en conservant cette double compétence : comprendre la technique et argumenter la décision.

À long terme, je souhaite développer mes propres projets technologiques, notamment Seshat, et contribuer à la croissance de l'entreprise familiale, Lyne \& Frères, active au Bénin dans l'agriculture, l'élevage et le commerce. Seshat est encore un projet entrepreneurial en développement. Ces deux ambitions demandent précisément ce que la maîtrise doit m'apporter : définir une stratégie de développement, gérer des ressources, évaluer des occasions et piloter une croissance. Mon temps à Montréal, lui, sera consacré aux études.

J'ajoute que cette étape doit aussi me permettre de travailler au Québec, dont j'ai besoin de comprendre le marché et les pratiques. Le cours Expérience Québec, réservé aux étudiants du cheminement avec projet supervisé, répond à cette attente.

\newpage
\section{3. Champs d'intérêt}
{\itshape\footnotesize Quels sont vos champs d'intérêt pour ce qui est de la recherche ou du projet supervisé?}\par
\vspace{0.35em}

Mon intérêt porte sur la mesure de la performance et le pilotage stratégique : la manière dont une organisation traduit une intention (croître, se diversifier, réduire ses coûts) en indicateurs, et la manière dont ces indicateurs reviennent ensuite orienter la décision.

À cette question s'en ajoute une seconde, que je veux travailler pendant la maîtrise : la finance d'entreprise. Je sais produire un indicateur ; je comprends moins bien comment une entreprise arbitre ses ressources, se finance et décide d'une acquisition. La spécialisation stratégie offre précisément le cours Finance d'entreprise et un cours sur les entreprises familiales : ce sont les deux que je veux suivre en priorité.

Cette question vient de mon travail. Dans l'assurance, j'ai passé six mois à fiabiliser des indicateurs de qualité des données et des indicateurs métier, et j'y ai vu qu'un chiffre contesté ne sert à personne : une part importante du travail consiste à établir la confiance dans la mesure avant même de pouvoir parler de ce qu'elle éclaire. Chez Lyne \& Frères, j'ai construit des tableaux de bord pour une direction qui n'avait pas d'analyste, et j'y ai mesuré l'écart entre un indicateur bien fait et un indicateur réellement utilisé.

Parmi les projets supervisés que propose HEC, la conception d'un tableau de bord pour mesurer l'atteinte des objectifs stratégiques correspond exactement au mandat que je recherche : il commence par la question de l'organisation (où veut-elle aller, sur quels leviers) et se termine par un instrument que les responsables doivent pouvoir s'approprier. La veille stratégique et l'analyse de l'avantage concurrentiel m'attirent également : j'ai eu à comparer et à arbitrer des choix technologiques pour mon propre produit sous contrainte de ressources, et l'analyse externe est un exercice que je n'ai jamais mené à l'échelle d'une entreprise.

Je privilégie une démarche appliquée, orientée vers la résolution de problèmes organisationnels concrets. Mon attente envers le projet supervisé est d'apprendre à cadrer un problème mal posé, à construire un diagnostic recevable et à défendre des recommandations devant des interlocuteurs qui n'ont pas d'intérêt a priori pour elles. C'est la partie du métier que je n'ai pas encore pratiquée, et c'est celle que je viens chercher.

\newpage
\section{4. Expérience de travail (facultatif)}
{\itshape\footnotesize Si vous possédez une expérience de travail pertinente pour votre admission, veuillez la décrire brièvement.}\par
\vspace{0.35em}

Mon expérience professionnelle tient en quatre étapes. Six mois chez Allianz France, ingénieur data : migration d'un datamart de l'assurance de SAS vers Python/PySpark, fiabilisation d'indicateurs de pilotage et mise en place de contrôles de qualité des données. Deux ans chez Lyne \& Frères, analyste de données : analyses commerciales, opérationnelles et agricoles, tableaux de bord et automatisations au service des décisions de la direction. Quatre mois chez Groupe Sylvagreg, en R\&D sur un sujet de vision par ordinateur. Depuis septembre 2026, je développe Seshat, mon propre projet de solutions d'intelligence artificielle et d'automatisation.

Ces expériences ont un point commun, qui est aussi ce que je viens consolider : chacune m'a demandé de traduire une question d'entreprise en travail analytique, puis d'en rendre le résultat utilisable par des personnes qui ne sont pas techniciennes. Mes expériences dans de petites structures et le développement de Seshat m'ont également confronté à la nécessité d'arbitrer entre plusieurs priorités avec des ressources limitées.

\vspace{0.9em}
Je me tiens à votre disposition pour tout renseignement complémentaire ou pour un entretien, et vous remercie de l'attention portée à ma candidature.

\vspace{1em}
{\bfseries Stéphane KPOVIESSI}

\end{document}
